\section{Introduction}
\label{sec:introduction}

For integers \(s,t\ge2\), the Ramsey number \(r(s,t)\) is the least
integer \(n\) such that every simple graph on \(n\) vertices contains
either a complete graph \(K_s\) or an independent set of size \(t\).
We determine the logarithmic exponent in the off-diagonal case \(s=5\).
We write \(\alpha(G)\) for the largest size of an independent set in \(G\).
All logarithms are natural.

\begin{theorem}
\label{thm:main}
There is an absolute constant \(C>0\) such that, for every
\(\varepsilon>0\), all sufficiently large integers \(t\) satisfy
\[
 \frac{t^4}{(\log t)^{3+\varepsilon}}
 \ \le\ r(5,t)\ \le\
 C\frac{t^4}{(\log t)^3}.
\]
The threshold for \(t\) may depend on \(\varepsilon\).
Consequently,
\[
 \lim_{t\to\infty}
 \frac{4\log t-\log r(5,t)}{\log\log t}=3.
\]
\end{theorem}

The theorem resolves the sharp logarithmic-exponent question for
\(r(5,t)\). It leaves open a constant-factor asymptotic for
\(t^4/(\log t)^3\).

\subsection{History and contribution}

Erd\H{o}s and Szekeres proved the classical bound
\(r(s,t)\le\binom{s+t-2}{s-1}\) \cite{ErdosSzekeres1935}.
Ajtai, Koml\'os and Szemer\'edi improved its fixed-\(s\) order to
\(O(t^{s-1}/(\log t)^{s-2})\) \cite{AKS1980}.
Li, Rousseau and Zang subsequently obtained leading upper constant
\(1+o(1)\) for fixed \(s\) \cite{LiRousseauZang2001}.
We include an elementary proof of the upper bound with an absolute
constant for \(s=5\), using uniform independent sets and random sampling.

The lower-bound problem has followed a different course.  Kim proved
that \(r(3,t)\) has order \(t^2/\log t\), matching the classical upper
bound in order \cite{Kim1995}.  For \(s=5\), Spencer's local-lemma
method gave a lower bound of order \((t/\log t)^3\)
\cite{Spencer1977}; the analysis of the random clique-free process by
Bohman and Keevash improved this to
\(t^3/(\log t)^{8/3}\), up to a constant factor
\cite[Theorem~1.2]{BohmanKeevash2010}.  Thus even the power of \(t\)
in the upper bound remained out of reach by these methods.
Finite geometry supplied a different route: Mattheus and Verstra\"ete
proved \(r(4,t)\ge c t^3/(\log t)^4\) using a construction based on
Hermitian unitals \cite{MattheusVerstraete2024}.

Brada\v{c}'s projective construction gives
\[
 r(s,t)\ge c_s\frac{t^{s-1}}{(\log t)^{2s-4}}
 \qquad(s\ge3,\ t\ge2)
\]
\cite[Theorem~1.1]{Bradac2026}.
For \(s=5\), its logarithmic exponent is \(6\).
We use the ordered projective-flag graph and adapt the marking
argument from Section~2.5 of that paper. The improvement here is in
the analysis of long independent sequences, closing the remaining
gap from exponent \(6\) to exponent \(3+o(1)\).
Earlier graphs on ordered incidence pairs appear in the work of
Codenotti, Pudl\'ak and Resta and of Kostochka, Pudl\'ak and R\"odl
\cite{CodenottiPudlakResta2000,KostochkaPudlakRodl2010}.
The distinction between steps with few choices and steps which
substantially reduce the remaining choices is also central to
Alon and R\"odl's independent-set counting method
\cite{AlonRodl2005}.  Our compression must account for the extra bias
created by choosing an independent sequence from a random stream.

The principal additional ingredient is a description theorem for
sparse point--hyperplane pairs. If two large sets have much fewer
incidences than the ambient density \(1/q\), a short message describes
a moderately larger set containing a constant fraction of the smaller
set. The message consists of samples selected from public random
tables, together with geometric data. A score formed from several
independent Poisson samples distinguishes most training points from
most ambient points. Polynomial bounds for rich lines and planes
control the dependencies in its high moments.

This description theorem is used inside an entropy argument.
The consistent sequence selected from a random stream has a
possibly highly biased law. We retain that law throughout, using
conditional entropy to choose representatives and to pay for every
message. A balanced tree of incidence tests successively reduces the
sets available to the decoder. This separates the geometric
description problem from the question of how to use its output
under adaptive conditioning.

\subsection{Proof overview}

Fix a small \(\eta>0\) and a large prime \(q\). We construct a
\(K_5\)-free graph on roughly \(q^4\log q\) positions whose independent
sets have fewer than \(q(\log q)^{1+\eta}\) positions.
A position receives a random incident pair \((a,b)\), where \(a\)
is a point of \(\PG(4,q)\) and \(b\) is a hyperplane containing it.
For positions \(i<j\), place an edge when \(a_i\) lies on \(b_j\)
but \(a_j\) does not lie on \(b_i\). A five-clique would give five
linearly independent points in one hyperplane.

Suppose every stream contained a long independent sequence, and
select one. We call its ordered flag tuple \emph{consistent}.
Even after auxiliary choices and conditioning on an event of
probability bounded below, each possible selected tuple is unlikely:
it must occur at some set of positions in the original stream.
This gives an entropy lower bound exceeding \(4\log q\) per flag
by a term of order \(\eta\log\log q\).

Two marking scans give almost every flag a set of at most \(Cq^4\)
joint possibilities. They also give separate bounds for its two
endpoints. Revealing a small number of representative flags lowers
the remaining mutual information. Consistency then has a geometric
consequence: most nearly independent pairs of endpoint laws have
very few incidences, except for collisions inside specific subspace
rectangles. A uniform occupancy bound on the original stream controls
those collisions.

The sparse-pair description theorem now applies to suitable nearly
uniform parts of the endpoint laws. Its output supplies samples
for fresh incidence tests. A balanced tree arranges these tests so
that the total description length is small. The tests retain a
constant fraction of the sequence while reducing the decoder's
available domains. A bounded number of repetitions makes the
marking description shorter than the entropy lower bound, a
contradiction.

Sections~\ref{sec:upper} and~\ref{sec:flags} establish the upper bound
and the flag-stream framework. The sparse-pair theorem is stated
in Section~\ref{sec:predictor-statement}; the following compression
argument proves the stream construction from it. Its proof is
then given in Sections~\ref{sec:predictor-setup}--\ref{sec:predictor-completion},
including the finite-field geometry and the score estimates.
Section~\ref{sec:conclusion} converts the stream parameters to
Theorem~\ref{thm:main}.
